%% ══════════════════════════════════════════════════════════════
%%  CHAPTER 9 — MEMORY & STATE SYSTEM
%% ══════════════════════════════════════════════════════════════
\chapter{Memory and State System}
\label{ch:memory}

\section{Overview}
\label{sec:memory:overview}

\autoscholar operates with a multi-layered state management system that spans the lifetime of individual sessions, pipeline runs, and the entire platform.
Unlike traditional agent systems with explicit long-term memory stores, \autoscholar relies primarily on the LLM's context window as its ``working memory'' and the database as its persistent state.

This chapter documents how state is managed, how previous work influences new papers, and what mechanisms exist for avoiding duplication.

\section{Context Window as Working Memory}
\label{sec:memory:context}

\subsection{The 1M Token Context}

Claude Opus 4.6's one-million-token context window serves as the primary working memory during paper generation.
Within a single pipeline run, all of the following accumulate in the context:

\begin{itemize}
  \item The system prompt (1,798 lines, $\sim$30K tokens).
  \item The research prompt (enhanced, 400--700 words).
  \item Literature search results (15--30 papers with abstracts and BibTeX).
  \item Data collection results (API responses, dataset summaries).
  \item Python code and execution outputs (scripts, figures, errors).
  \item Written paper sections (\LaTeX{} content).
  \item Tool results and governor directives.
  \item Conversation history (all previous turns).
\end{itemize}

\subsection{Token Budget Management}

The system manages the context budget through two mechanisms:

\begin{description}
  \item[Tool result truncation.] Results exceeding 15,000 characters are truncated with a ``[truncated]'' marker. This is applied selectively---larger results (e.g., full API response dumps) are truncated while smaller results are preserved in full.

  \item[History summarization.] When estimated message tokens exceed the 500K safety threshold (half the 1M window), the system calls Claude to compress older conversation history:
  \begin{enumerate}
    \item The first 50\% of messages are selected for summarization.
    \item Claude generates a concise summary preserving key findings, data sources, and methodology decisions.
    \item The original messages are replaced with a single summary message.
    \item This frees context space for the remaining pipeline stages.
  \end{enumerate}
\end{description}

\subsection{Implications}

Within a single paper generation run, the system has excellent ``memory''---it can reference any earlier finding, dataset, or decision because all information remains in context.
However, there is no memory \emph{across} runs: each paper starts with a fresh context, and the system cannot learn from previous papers' successes or failures.

\section{Session State}
\label{sec:memory:session}

\subsection{Conversation Sessions}

Each interactive chat or pipeline run is tracked as a conversation session:

\begin{description}
  \item[\texttt{conversationSessions} table:] Stores session metadata including user ID, session ID, title, message count, token usage, cost, and timestamps.
  \item[\texttt{messages} table:] Stores individual messages with role (user/assistant), content, serialized sources, and timestamps.
\end{description}

\subsection{Real-Time State Broadcasting}

The SessionBroadcaster (\file{sessionBroadcaster.ts}, 180 lines) manages real-time state for live monitoring:

\begin{itemize}
  \item Maintains per-session event buffers (last 200 events).
  \item Manages subscriber sets for SSE streaming.
  \item Tracks session activity with last-heartbeat timestamps.
  \item Supports reconnection with buffered event replay (catch-up for disconnected clients).
  \item Automatic cleanup of sessions idle for 30 minutes or older than 1 hour.
\end{itemize}

\subsection{Session Persistence}

Session data is persisted every 8 seconds during active generation:
\begin{itemize}
  \item Messages (question-answer pairs).
  \item Tool results (structured JSON).
  \item Figures (base64-encoded or URL references).
  \item Token counts and cost accumulation.
  \item Agent step progress.
\end{itemize}

\section{Pipeline State}
\label{sec:memory:pipeline}

\subsection{Governor State}

The \governor maintains typed state throughout a pipeline run (see \cref{sec:agents:governor}).
This state includes:

\begin{itemize}
  \item Current pipeline stage.
  \item Active prompt profile.
  \item Error counts (Python, \LaTeX{}, data).
  \item Retry budgets (3/3/2 for Python/\LaTeX{}/data).
  \item Confidence score (0.0--1.0).
  \item Figure and word counts.
  \item Warning accumulator.
  \item Publication verdict.
\end{itemize}

Governor state is serialized to JSON and stored in the \texttt{paper\_runs.governor\_state} column, enabling:
\begin{itemize}
  \item Post-mortem analysis of failed runs.
  \item Checkpoint resume with full governor state recovery.
  \item Historical analysis of governor intervention patterns.
\end{itemize}

\subsection{Decision Log}

All governor decisions are recorded in the DecisionLog with:
\begin{itemize}
  \item UUID per decision.
  \item Timestamp.
  \item Pipeline stage.
  \item Trigger description.
  \item Action taken.
  \item Reasoning.
  \item Confidence delta.
  \item Profile transition (if any).
\end{itemize}

The decision log persists in both the \texttt{paper\_runs} table (serialized JSON) and the \texttt{governor\_decisions} table (individual records).

\subsection{Checkpoint State}

The checkpoint system captures a snapshot of the entire pipeline state:

\begin{itemize}
  \item Full conversation history (all Claude messages).
  \item Governor state and decision log.
  \item Tool results accumulated so far.
  \item Written paper sections.
  \item Figure count and metadata.
  \item Pipeline stage at checkpoint time.
\end{itemize}

Checkpoints enable crash recovery: after a server restart, an administrator can resume a run from its last checkpoint.

\section{Cross-Paper Memory}
\label{sec:memory:cross-paper}

\subsection{Current State: No Cross-Paper Memory}

\autoscholar currently has \textbf{no mechanism for cross-paper memory}.
Each paper generation run starts with a fresh context window, and the system has no awareness of:

\begin{itemize}
  \item Topics previously covered.
  \item Methods previously applied.
  \item Data sources previously used.
  \item Findings from previous papers.
  \item Reviewer feedback from previous runs.
\end{itemize}

\subsection{Implications}

This absence has several consequences:

\begin{description}
  \item[Topic duplication.] While the sequential prompt selection ensures topic diversity, there is no mechanism to prevent the system from revisiting a topic if the prompt bank contains similar entries or if prompts are selected non-sequentially.

  \item[No learning from failures.] If a particular data source consistently fails or a method doesn't converge for certain data types, the system cannot learn to avoid these combinations.

  \item[No citation of own work.] The system cannot cite its own previously published papers, even when they would be relevant.

  \item[No methodological improvement.] The system cannot improve its analysis techniques based on reviewer feedback from previous runs.
\end{description}

\subsection{Implicit Memory via Prompt Bank}

The prompt bank provides a form of implicit memory:
\begin{itemize}
  \item Sequential selection ensures each prompt is used exactly once before cycling.
  \item The domain-method compatibility matrix prevents nonsensical combinations.
  \item Topic diversity is built into the bank's construction (14 topics per domain, interleaved iteration order).
\end{itemize}

However, this is ``memory'' in the weakest sense---it is pre-programmed diversity, not learned adaptation.

\section{Platform State}
\label{sec:memory:platform}

\subsection{Site Configuration}

The \texttt{site\_config} table provides a key-value store for runtime configuration:

\begin{itemize}
  \item Scheduler settings (cadence, parallel limit, thresholds).
  \item Custom system prompt overrides.
  \item Feature flags.
  \item Cost tracking baselines.
\end{itemize}

Values persist across server restarts and are modifiable via the admin API.

\subsection{Analytics Accumulation}

Platform-wide analytics accumulate over time:

\begin{itemize}
  \item Total papers generated (published + rejected).
  \item Total tokens consumed and API cost.
  \item Per-domain paper counts and quality scores.
  \item Monthly publication timelines.
  \item Python execution counts and code line totals.
\end{itemize}

\subsection{Published Paper Corpus}

The published paper corpus serves as the platform's ``institutional memory'':

\begin{itemize}
  \item Full \LaTeX{} source and BibTeX for every published paper.
  \item Peer review scores and feedback.
  \item Domain classification and methodology tags.
  \item Reproducible project archives.
\end{itemize}

While this corpus is not currently used to inform new paper generation, it represents a rich dataset that could power future cross-paper learning (see \cref{ch:improvements}).

\section{State Consistency and Recovery}
\label{sec:memory:consistency}

\subsection{Crash Recovery}

State consistency after crashes is maintained through:

\begin{enumerate}
  \item \textbf{8-second checkpoint saves:} Session state is persisted frequently enough to minimize data loss.
  \item \textbf{Heartbeat monitoring:} Stalled runs are detected within 5 minutes and marked as failed.
  \item \textbf{Force-fail on restart:} All active runs are force-failed on server restart to prevent zombie processes.
  \item \textbf{Manual resume:} Administrators can resume runs from checkpoints via the admin API.
\end{enumerate}

\subsection{Data Integrity}

\begin{itemize}
  \item SQLite uses WAL (Write-Ahead Logging) mode for crash-safe writes.
  \item PostgreSQL provides full ACID guarantees.
  \item ASID generation uses atomic database operations to prevent collision under concurrent access.
  \item File operations use temporary directories with rename-on-success to prevent partial writes.
\end{itemize}
